\section{Accessible tangent geometry}
\subsection{From quantum sensitivity to accessible sensitivity}
Let $|\psi_{\bm\theta}\rangle$ be a differentiable pure state and $X_v$ the horizontal tangent generated by parameter-space direction $v$. Its directional QFI is $F_Q(v)=4\langle X_v,X_v\rangle$. A measurement $\mathcal M$ maps the quantum tangent to a classical score tangent, and a retained readout space $\mathcal A$ projects that score onto directions represented by the chosen observables. The spectral accessibility framework packages the two contractions into
\begin{equation}
 T_{\mathcal M,\mathcal A}=A_{\mathcal M}^{\dagger}P_{\mathcal A}A_{\mathcal M},\qquad 0\preceq T_{\mathcal M,\mathcal A}\preceq I,
\end{equation}
so that for a normalized quantum tangent $u$,
\begin{equation}
 \frac{I_{\mathcal M,\mathcal A}(u)}{F_Q(u)}=\langle u,T_{\mathcal M,\mathcal A}u\rangle.
\end{equation}
For a differentiable classical head built from retained expectations, the effective first-order observable remains in the retained span. The first-order loss response is therefore bounded by accessible information, motivating the use of the same tangent span to define the optimizer metric.

\subsection{Accessible pullback metric}
The full QNG metric is the pure-state QFIM pullback
\begin{equation}
 (\GQ)_{ij}=4\,\mathrm{Re}\!\left[\langle\partial_i\psi|\partial_j\psi\rangle-\langle\partial_i\psi|\psi\rangle\langle\psi|\partial_j\psi\rangle\right].
\end{equation}
Let $\Phi_{\mathcal M,\mathcal A}$ map a horizontal quantum tangent to its retained covariance-normalized readout score. We define
\begin{equation}
 (\GA)_{ij}=\left\langle\Phi_{\mathcal M,\mathcal A}(\partial_i\psi),\Phi_{\mathcal M,\mathcal A}(\partial_j\psi)\right\rangle.
 \label{eq:GA}
\end{equation}
Thus $\GA$ is the parameter-space pullback of accessible tangent geometry and ideally satisfies $0\preceq\GA\preceq\GQ$. The damped updates compared here are
\begin{align}
 \Delta\bm\theta_{\rm QNG}&=-\eta(\GQ+\lambda I)^{-1}\nabla L,\\
 \Delta\bm\theta_{\rm AQNG}&=-\eta(\GA+\lambda I)^{-1}\nabla L.
\end{align}
The goal is not $\GA\simeq\GQ$ in matrix norm. The operational question is whether the accessible preconditioner preserves or improves directions that matter to the task.

\subsection{Nested readout probes}
The implementation diagnoses three nested diagonal readout spaces: a minimal probe $A_1$, the full local-$Z$ span, and the joint span of diagonal Pauli strings through weight two, denoted $\mathcal A_{\le2}$. The main AQNG metric uses the $\le2$ accessible geometry. The expected nesting is
\begin{equation}
 G_{A_1}\preceq G_{\rm local}\preceq G_{\le2}\preceq\GQ.
\end{equation}
We record metric rank, trace retention $\Tr(G_{\mathcal A})/\Tr(\GQ)$, and cosine between the resulting preconditioned direction and the full-QNG direction.
